% chapter_01_ME.tex -- Chapter 01, ME: the matrix element q qbar -> b bbar of the ONE event, and nothing else.
% Script: test/01_ME/run_01_ME.sh; numbers from test/01_ME/out/ (step1_GEN_01_ME.log, genPartAnalyzer_01_ME.log,
% ledger_01_ME.txt); every number is listed in full in doc/chapter/chapter_01_ME.md.
\chapter{01 -- ME: two partons in, a \bbbar{} pair out, and nothing else}
\label{ch:01_ME}

\section{What this chapter is about}

Every chapter of this book follows the same proton--proton collision. This first chapter produces it in its
barest form: Pythia8 is asked to pick two partons out of the two protons, to let them scatter into a
$b$ quark and a $\bar b$ antiquark through the lowest-order matrix element of $\qqbar\to\bbbar$, and to stop
there. The initial- and final-state parton showers, the multiparton interactions, the hadronisation and
the hadron decays are all switched off (\code{PartonLevel:ISR/FSR/MPI = off}, \code{HadronLevel:all = off}),
so that the output file holds the hard process, the two incoming partons after a small transverse kick,
and the beam remnants: fourteen entries in \code{genParticles}, six of them final.\footnote{All numbers of
this chapter are read from \path{test/01_ME/out/step1_GEN_01_ME.log} (the Pythia8 listing of the event),
\path{genPartAnalyzer_01_ME.log} and \code{ledger\_01\_ME.txt}; the companion
\path{doc/chapter/chapter_01_ME.md} lists them in full.}

The chapter explains what a matrix element is, how the two momentum fractions $x_1$ and $x_2$ are drawn
from the parton distribution functions, why the event comes out so lopsided (a $b$ at $\eta=-1.44$ with
$252\GeV$, a $\bar b$ at $\eta=-3.75$ with $2416\GeV$), how to read the Pythia8 status codes of the
record, and what the \code{GEN} step of \code{cmsDriver.py} does around Pythia8: the vertex smearing, the
conversion to \code{genParticles}, and the generator jets and missing momentum computed already here.

\begin{plainwords}[A matrix element]
A matrix element is the quantum-mechanical amplitude for one definite process at the parton level, here
$u\bar u\to b\bar b$ through a gluon, computed to a fixed order of perturbation theory (here the lowest
one, $\alpha_s^2$). Its square, integrated over the phase space of the outgoing partons and folded with
the probability of finding the incoming partons in the protons, is the cross section. The matrix element
knows nothing about radiation, about the other partons of the protons or about hadrons: it is the one
$2\to2$ scattering in the middle of the event, which the next chapters dress but do not change.
\end{plainwords}

\section{From the previous chapter}

There is no previous chapter: this is the starting point of the journey, so this section describes the
inputs. The physics is defined by one generator fragment,
\path{Configuration/GenProduction/python/BBbar_14TeV_TuneCP5_cfi.py} of the package (a symlink is made where
\code{cmsDriver.py} looks for it). A fragment is a small python file that configures the \code{generator}
module (a \code{Pythia8ConcurrentGeneratorFilter}): it switches on the single process
\code{HardQCD:qqbar2bbbar} and switches off its gluon-initiated partner \code{HardQCD:gg2bbbar}, loads the
two parameter blocks that every CMS Pythia8 sample uses (\code{pythia8CommonSettings} from
\path{Configuration/Generator/python/Pythia8CommonSettings_cfi.py} and the tune \code{pythia8CP5Settings}
from \path{Configuration/Generator/python/MCTunes2017/PythiaCP5Settings_cfi.py}, which also selects the
PDF set \code{NNPDF31\_nnlo\_as\_0118}), and asks for a minimum transverse momentum of the hard scattering
\code{PhaseSpace:pTHatMin = 30}. The script \path{test/01_ME/run_01_ME.sh} (through \code{gen\_stage} of
\path{test/common/jitjet_common.sh}) overrides two things with \code{-{}-customise\_commands}: the
centre-of-mass energy, $14\TeV$ in the fragment, becomes $13.6\TeV$ (the 2024 LHC), and \ptHat{} is
raised to $100\GeV$, so that the two $b$ quarks make jets that the later chapters can follow.

Two more inputs make the event reproducible. The random seeds are the defaults of
\path{IOMC/RandomEngine/python/IOMC_cff.py} (\code{generator} 123456789, \code{VtxSmeared} 98765432) and
are never changed, and exactly one event is generated (\code{-n 1}: run 1, lumi 1, event 1). Pythia8
draws the hard process before anything else, so chapters 02 to 06, which rerun the same fragment with the
same seeds and one more step switched on, obtain the same $\qqbar\to\bbbar$ kinematics. The next
chapter, 02\_PS, switches on the showers and dresses the four partons of this record with radiation.

\section{The numbers of the event}

\paragraph{The hard process.} Pythia8 prints the event it generated in the \code{cmsRun} log
(\code{step1\_GEN\_01\_ME.log}, "PYTHIA Info Listing" and "Event Listing (hard process)");
table~\ref{tab:01_hard} collects it. The two protons of $6800\GeV$ come in along $\pm z$. Beam~A (the
proton flying towards $+z$) supplies a $\bar u$ sea antiquark with the momentum fraction
$x_1 = 2.17\times10^{-3}$, that is $14.8\GeV$; beam~B supplies a $u$ valence quark with $x_2 = 0.390$,
that is $2653\GeV$. The parton distribution functions at the scale of the scattering,
$Q^2 = 1.284\times10^4\GeV^2$ ($Q = 113.3\GeV$, the transverse mass of the $b$), are
$x f_{\bar u}(x_1) = 1.145$ and $x f_u(x_2) = 0.211$: the sea is dense at small $x$, the valence $u$ is
still sizeable at $x=0.39$. Their invariant mass is $\sqrt{\shat} = 396.1\GeV$, the strong coupling at $Q$ is
$\alpha_s = 0.114$. Pythia8's estimate of the cross section after this single event, $157.7\pm157.7$ pb
(\code{GenXsecAnalyzer} at the end of the log), is not a measurement; the note below the table explains why
and gives $348.6\pm6.0$ pb from 1000 events.

\begin{table}[htb]
\centering\small
\caption{The hard process of the event, from the Pythia8 listing in \code{step1\_GEN\_01\_ME.log}
(momenta in GeV; the status codes are Pythia8's: 4 beam, 21 incoming, 23 outgoing).}
\label{tab:01_hard}
\begin{tabular}{llrrrrrrrr}
\toprule
 & particle & status & $x$ & $p_x$ & $p_y$ & $p_z$ & $E$ & \pt & $\eta$ \\
\midrule
beam A & $p$ & 4 & & 0 & 0 & $+6800.0$ & 6800.0 & 0 & \\
beam B & $p$ & 4 & & 0 & 0 & $-6800.0$ & 6800.0 & 0 & \\
in 1 & $\bar u$ (sea) & 21 & 0.00217 & 0 & 0 & $+14.785$ & 14.785 & 0 & \\
in 2 & $u$ (valence) & 21 & 0.390 & 0 & 0 & $-2652.934$ & 2652.934 & 0 & \\
out & $\bar b$ & 23 & & 62.071 & 94.690 & $-2413.001$ & 2415.661 & 113.22 & $-3.75$ \\
out & $b$ & 23 & & $-62.071$ & $-94.690$ & $-225.148$ & 252.059 & 113.22 & $-1.44$ \\
\bottomrule
\end{tabular}
\end{table}

\paragraph{Note: a cross section needs many events, not one.} Pythia8 computes the cross section by Monte
Carlo integration. It throws trial points in the phase space of $\qqbar\to\bbbar$ with $\ptHat>100\GeV$, gives
each the weight matrix element $\times$ the two PDFs, keeps the point as an event with the probability weight /
maximum (the maximum, $1.004\times10^{-6}\mb = 1004$ pb, is found at initialisation), and quotes the average weight
of all trials as $\sigma$. After the single event of the book there was one trial: the estimate is the weight
of that one point, $157.7$ pb (it was kept with the probability $157.7/1004 = 16\%$), and with one value there
is no spread to take an error from, so Pythia8 quotes $\pm157.7$ pb, a 100\% uncertainty. The error of an average
falls like $1/\sqrt{N}$, so \path{run_01_ME.sh} runs a second GEN job with the same fragment, settings and seeds and
1000 events, written to no file (\code{gen\_xsec} of \path{jitjet_common.sh}, log \path{xsec_GEN_01_ME.log}).
Its first event is the event of the book, and the 1000 events came from 2913 trials:
\[
  \sigma(\qqbar\to\bbbar,\ \ptHat>100\GeV,\ \sqrt{s}=13.6\TeV) = 348.6\pm6.0~\mathrm{pb}\quad(1.7\%).
\]
A quick check: the fraction of trials kept times the maximum, $1000/2913\times1004 = 345$ pb, agrees within the
error. The single-event value is 2.2 times too small: one weight says where the event sits on the matrix element,
not what the total is. This is the cross section of the one subprocess the fragment switches on, not of all \bbbar{}
production. It is the number to quote; \code{XSEC\_EVENTS} changes the number of events.

\paragraph{Why the event is so asymmetric.} The two $b$ quarks have the same \pt{} of $113.2\GeV$ and
are back to back in $\phi$ ($\Delta\phi = 3.14$, \path{genPartAnalyzer_01_ME.log}), as they must be for a
$2\to2$ process with no radiation. Everything else about them is different: the $b$ carries $252\GeV$
at $\eta = -1.44$, the $\bar b$ carries $2416\GeV$ at $\eta = -3.75$. Two things add up. First, the
parton--parton system is not at rest in the laboratory: its rapidity is
$y = \tfrac12\ln(x_1/x_2) = -2.595$, so the whole \bbbar{} pair is boosted towards $-z$, the direction of
the valence $u$ quark that brought almost all the energy. Second, in the rest frame of the two partons the
scattering angle is $\hat\theta = 145^\circ$ (Pythia8 quotes it for the $\bar b$), which gives the $\bar b$
a rapidity of $-1.157$ in that frame and the $b$ one of $+1.157$. Adding the two,
$-2.595 - 1.157 = -3.752$ and $-2.595 + 1.157 = -1.438$, reproduces the two pseudorapidities of the table (the
$b$ mass makes rapidity and pseudorapidity differ by only $0.001$). The energies follow from $E \simeq \pt\cosh\eta$:
$113.2\cosh(1.44) = 252\GeV$ and $113.2\cosh(3.753) = 2415\GeV$. In the 2024 detector the $b$ points at the barrel--endcap transition of the calorimeters and the
$\bar b$ at the forward hadron calorimeter, outside the tracker: two identical quarks that the detector
will see in two very different ways.

\begin{plainwords}[Sampling the protons]
A proton is not a particle with a fixed content: each time it is probed at a scale $Q$, the probability
to find a parton of type $i$ with a fraction $x$ of its momentum is $f_i(x, Q^2)$, the parton distribution
function, measured in other experiments and tabulated (here NNPDF3.1 through LHAPDF, see the
\code{LHAPDF 6.4.0 loading} lines of the log). Pythia8 draws $x_1$ and $x_2$ from these functions and the
scattering angle from the matrix element, keeps the event if $\ptHat > 100\GeV$, and tries again
otherwise. Most draws pair a small $x$ with a large one, because the small-$x$ sea is so much denser:
lopsided events like this one are the rule, not the exception.
\end{plainwords}

\paragraph{The record and its status codes.} Table~\ref{tab:01_record} lists the fourteen entries that
\code{genParticles} holds after this stage, in the order of the file. CMSSW keeps the Pythia8 status
code of each entry (Pythia8 manual, "Particle Properties"): 4 is a beam proton, 21 and 23 are the
incoming and outgoing partons of the hardest subprocess, 61 is an incoming parton with its primordial
$k_{\mathrm T}$ included, 1 is a final-state particle. Codes 31--39 (multiparton interactions), 41--59
(showers) and 71--89 (hadronisation) appear in the later chapters; 63, the outgoing beam remnant, is
written as status 1 because nothing else happens to it here. The record is a history: a particle appears
once per copy, linked to its mothers and daughters. The two incoming partons (status 21, $\pt = 0$) are
copied with status 61 and a small transverse momentum, $1.37$ and $0.97\GeV$, the primordial
$k_{\mathrm T}$ (\code{BeamRemnants:primordialKThard}, default width $1.8\GeV$, not touched by the tune);
the same kick, $1.93\GeV$ in total, is passed to the \bbbar{} pair, so the final $b$ and $\bar b$
(status 1, $\pt = 113.56$ and $112.33\GeV$) no longer balance exactly. The rest of the protons goes into
the beam remnants: a $ud_1$ diquark with $6643\GeV$ and two $u$ quarks ($133$ and $9\GeV$) from beam~A, a
$ud_0$ diquark with $4147\GeV$ from beam~B, all close to the beam axis ($|\eta| > 3.3$).

\begin{table}[htb]
\centering\footnotesize
\caption{The fourteen entries of \code{genParticles} in \code{step1\_GEN\_01\_ME.root} (FWLite, see the
companion), with the Pythia8 status code, the energy, the transverse momentum and the pseudorapidity.
The six status-1 entries are the final state of this stage; their energies add up to $13600.0\GeV$.}
\label{tab:01_record}
\begin{tabular}{rllrrrl}
\toprule
\# & particle & status & $E$ [GeV] & \pt{} [GeV] & $\eta$ & role \\
\midrule
0 & $p$ & 4 & 6800.00 & 0 & $+\infty$ & beam A ($+z$) \\
1 & $p$ & 4 & 6800.00 & 0 & $-\infty$ & beam B ($-z$) \\
2 & $\bar u$ & 21 & 14.79 & 0 & & incoming, $x_1 = 0.00217$ \\
3 & $u$ & 21 & 2652.93 & 0 & & incoming, $x_2 = 0.390$ \\
4 & $\bar b$ & 23 & 2415.66 & 113.22 & $-3.75$ & outgoing, matrix element \\
5 & $b$ & 23 & 252.06 & 113.22 & $-1.44$ & outgoing, matrix element \\
6 & $\bar u$ & 61 & 14.82 & 1.37 & $+3.07$ & incoming with primordial $k_{\mathrm T}$ \\
7 & $u$ & 61 & 2652.93 & 0.97 & $-8.61$ & incoming with primordial $k_{\mathrm T}$ \\
8 & $\bar b$ & 1 & 2414.61 & 112.33 & $-3.76$ & final $\bar b$ \\
9 & $b$ & 1 & 253.15 & 113.56 & $-1.44$ & final $b$ \\
10 & $ud_1$ (2103) & 1 & 6643.29 & 0.74 & $+9.80$ & beam remnant of A \\
11 & $u$ & 1 & 132.84 & 0.41 & $+6.47$ & beam remnant of A \\
12 & $u$ & 1 & 9.10 & 0.63 & $+3.36$ & beam remnant of A \\
13 & $ud_0$ (2101) & 1 & 4147.02 & 0.97 & $-9.06$ & beam remnant of B \\
\bottomrule
\end{tabular}
\end{table}

\paragraph{What the GEN step does besides Pythia8.} The configuration written by \code{cmsDriver.py}
(\path{test/01_ME/step1_GEN_01_ME_cfg.py}) runs the path \code{generation\_step}: the generator, then the
sequence \code{pgen} of \path{Configuration/StandardSequences/python/Generator_cff.py}, which holds much
more than Pythia8:
\begin{itemize}
\item \code{VtxSmeared}: Pythia8 puts the collision at the origin; this module moves it to a point drawn
  from the 2024 luminous region. It is chosen by the option \code{-{}-beamspot DBrealistic}: the
  \code{BetafuncEvtVtxGenerator} configured in the file
  \path{IOMC/EventVertexGenerators/python/VtxSmearedRealistic_cfi.py} with \code{readDB = True}, which
  reads the beam parameters from the global tag \code{150X\_mcRun3\_2024\_realistic\_v2} (payload
  \code{SimBeamSpot\_Realistic25ns13p6TeV\allowbreak Early2024Collision\allowbreak VtxSmearing\allowbreak Parameters\_v1\_mc}):
  the centre $(0.0285, -0.0166, 0.416)\cm$, $\sigma_z = 3.6\cm$, $\beta^* = 30\cm$ and the emittance
  $5.725\times10^{-8}\cm$, that is $\sigma_{x,y} = 13\um$. This event lands at
  $(0.0277, -0.0169, 3.112)\cm$ and at $t = 0.121\ns$ (the \code{genParticles:t0} product); the time is a
  Gaussian of width $\sigma_z/c = 120\ps$ drawn independently of $z$. Only the vertex moves, no momentum
  changes. \code{generatorSmeared} is a copy of the result under the label everything downstream reads.
\item \code{genParticles} (\code{GenParticleProducer}): the HepMC record as the \code{reco::GenParticle}
  collection of table~\ref{tab:01_record}, with mother and daughter links.
\item the generator jets and missing momentum of \code{genJetMET}: \code{genParticlesForJets} selects the
  status-1 particles, and \code{ak4GenJets}, \code{ak8GenJets} (anti-$k_{\mathrm T}$, $R = 0.4$ and $0.8$,
  \code{jetPtMin = 3}, \path{RecoJets/JetProducers/python/GenJetParameters_cfi.py}) and their \code{NoNu}
  variants are already run here. Here each collection holds two jets of one constituent each, the $b$ ($\pt = 113.56\GeV$) and the
  $\bar b$ ($112.33\GeV$); the remnants are below threshold. \code{genMetTrue} is $5\times10^{-14}\GeV$
  (no neutrinos, momentum conserved); \code{genMetCalo}, which only counts $|\eta| < 5$, is $1.40\GeV$.
\item \code{LHCTransport} (\code{PPSSimTrackProducer}, Run-3 era): the transport of intact protons to
  the PPS Roman pots; this event has none.
\end{itemize}

\section{Ledger and display}

The ledger of this stage (table~\ref{tab:01_ledger}) compares the two protons that enter with the six
final-state entries that leave. Nothing is lost: the generator conserves four-momentum exactly. What the
matrix element did is to move $2677\GeV$ out of the beam direction into the region a detector can see
($|\eta| < 5.2$), almost all of it carried by the $\bar b$ in the forward calorimeter region; the
$10923\GeV$ of the diquarks and of the $u$ quark at $\eta = 6.47$ leave along the beam pipe and will be
dropped at the hand-over to Geant4 (chapter 09). The scalar \pt{} of the final state, $228.6\GeV$, is the
\pt{} of the two $b$ quarks plus the $2.7\GeV$ of the remnants.

\begin{table}[htb]
\centering\small
\caption{Ledger of the matrix-element stage: the two protons before, the six final-state entries
(status 1) after, from \code{ledger\_01\_ME.txt} ("by region" and "visible" lines). The regions are those of
the 2024 detector: tracker $|\eta| < 2.5$, forward calorimeter $3.0 < |\eta| < 5.2$, beyond $5.5$ nothing
is simulated; nothing sits at $2.5 < |\eta| < 3.0$ and there are no neutrinos.}
\label{tab:01_ledger}
\begin{tabular}{lrrr}
\toprule
quantity & before (two protons) & after (ME record) & difference \\
\midrule
particles (status 1) & 2 & 6 & $+4$ \\
total energy [GeV] & 13600.0 & 13600.0 & 0 \\
scalar \pt{} [GeV] & 0 & 228.6 & $+228.6$ \\
energy at $|\eta| < 2.5$ [GeV] (particles) & 0 (0) & 253.1 (1: the $b$) & $+253.1$ \\
energy at $3.0 < |\eta| < 5.2$ [GeV] (particles) & 0 (0) & 2423.7 (2: $\bar b$, $u$) & $+2423.7$ \\
energy at $|\eta| > 5.5$ [GeV] (particles) & 13600.0 (2) & 10923.1 (3) & $-2676.9$ \\
visible energy, $|\eta| < 5.2$ [GeV] & 0 & 2676.9 & $+2676.9$ \\
visible scalar \pt{}, $|\eta| < 5.2$ [GeV] & 0 & 226.5 & $+226.5$ \\
\midrule
hard process: $\sqrt{\shat}$ [GeV] & & 396.1 & \\
hard process: $x_1$, $x_2$ & & 0.0022, 0.390 & \\
\bottomrule
\end{tabular}
\end{table}

Figure~\ref{fig:01_display} is the $\eta$--$\phi$ display of the event at this stage, the first page of
the \code{GenPartAnalyzer} output: the two outgoing quarks of the matrix element, the $b$ in red and the
$\bar b$ in blue, the marker size growing with \pt{}; the incoming partons and the beam remnants are
along the beam axis and are not drawn. The two markers sit at the same $\pt$, opposite in $\phi$, $2.3$
units of $\eta$ apart. The following chapters redraw this page with more and more particles around them.

\begin{figure}[htb]
\centering
\includegraphics[width=0.84\textwidth]{01_ME/out/genPartAnalyzer_01_ME_page-1.png}
\caption{The event after the matrix element: the outgoing $b$ (red) and $\bar b$ (blue) of the hard
process in the $\eta$--$\phi$ plane, with their \pt{} (\code{test/01\_ME/out/genPartAnalyzer\_01\_ME\_page-1.png}).
The right panel lists the Pythia8 switches of the stage; stages 2 to 5 do not exist in this file.}
\label{fig:01_display}
\end{figure}

\FloatBarrier
\section{Open problems for Phase-2}

The generator stage does not know which detector follows it, which is exactly why its settings are the
easiest to leave untouched. A look into \code{\$CMSSW\_RELEASE\_BASE/src} of CMSSW\_15\_0\_5 gives this list.

\paragraph{The beam energy: 14 TeV in the fragment, 13.6 TeV here, and a 13 TeV luminous region in the
Phase-2 workflows.} The fragment is modelled on \path{Configuration/Generator/python/TTbar_14TeV_TuneCP5_cfi.py}
and carries \code{comEnergy = 14000}; the script overrides it to $13600$ (\code{Beams:eCM} in the log). The Phase-2 relval workflows of the release use the 14 TeV fragments with the
scenario \code{GenSimHLBeamSpot}, defined for every \code{Run4} key in
\path{Configuration/PyReleaseValidation/python/upgradeWorkflowComponents.py}; the key \code{Run4D110},
for instance, sets the geometry to \code{ExtendedRun4D110}, the global tag to
\code{auto:phase2\_realistic\_T33} and the era to \code{Phase2C17I13M9}. That scenario (\code{upgradeStepDict['GenSimHLBeamSpot']} in \path{relval_steps.py})
runs \code{-{}-beamspot DBrealisticHLLHC} with the conditions \code{gt + '\_13TeV'}, and
\path{Configuration/AlCa/python/autoCondModifiers.py} (\code{autoCondBSHLLHC13TeV}) resolves the suffix by
appending the payload \path{SimBeamSpotHLLHCObjects_z4p3cm_t193ns_13TeV_nominal_mc_v1}, whose proton
energy is $6500\GeV$, to a job that generates $7000\GeV$ beams: the transverse size of the luminous
region, which scales as $1/\sqrt{\beta\gamma}$, comes out 4\% too wide. The 14 TeV payload
(\path{..._14TeV_nominal_mc_v1}, $E_p = 7000\GeV$) is only used by the variant \code{GenSimHLBeamSpot14},
which no \code{Run4} key runs by default.

\paragraph{The tune and its PDF are from 2017 and 13 TeV.} \code{pythia8CP5Settings}
(\path{MCTunes2017/PythiaCP5Settings_cfi.py}) is the CP5 tune of 2017 on \code{NNPDF31\_nnlo\_as\_0118},
$\alpha_s = 0.118$ at second order, with fixed total and elastic cross sections
\code{SigmaTotal:sigmaTot = 100.309}, \code{sigmaEl = 21.89} (\code{SigmaTotal:mode = 0}). The release
holds a Run-3 copy, \path{MCTunesRun3ECM13p6TeV/PythiaCP5Settings_cfi.py}, which differs in exactly two
numbers, \code{sigmaTot = 101.037} and \code{sigmaEl = 22.08}, the 13.6 TeV values; the fragment imports
the 2017 file, and there is no 14 TeV directory at all. For the matrix element (PDF, $\alpha_s$, process
switches) it makes no difference; it matters from chapter 04 on (the multiparton interactions use
\code{sigmaTot}) and for the minimum-bias library of chapter 08. The one energy-dependent tune parameter, $p_{\mathrm T0} = 1.41\,(\sqrt s/7\TeV)^{0.03344}\GeV$,
moves from $1.442$ to $1.443\GeV$ between 13.6 and 14 TeV. No newer tune or PDF set is wired into
\path{Configuration/Generator} of this release.

\paragraph{The \ptHat{} choice and the process selection are ours, not the release's.} The fragment asks
for $\ptHat > 30\GeV$ ("so that the b quarks produce jets reconstructable for JME studies", its own
comment) and switches \code{HardQCD:gg2bbbar} off; the script raises the cut to $100\GeV$. Both are
choices of the book: the gluon channel, which dominates inclusive \bbbar{} production, is excluded to
keep one clean $2\to2$ history, and \ptHat{} is a cut on the hard scattering before any radiation. At
14 TeV and $\mathrm{PU} = 200$ the Phase-2 trigger thresholds decide what \ptHat{} is useful; nothing in
the release ties the two together.

\paragraph{The vertex smearing is a GEN-step decision, and the time of the vertex is where 2024 and
HL-LHC differ most.} \code{cmsDriver.py} requires \code{-{}-beamspot} at the GEN step
(\path{Configuration/Applications/python/ConfigBuilder.py}, "GEN step always requires to have a
VtxSmearing scenario") and ignores it in a SIM step: a Phase-2 chain cannot re-use this chapter's file,
it has to regenerate. The two models are different programs:
\begin{itemize}
\item 2024, \code{DBrealistic}: \code{BetafuncEvtVtxGenerator} with the payload quoted above; in
  \path{IOMC/EventVertexGenerators/src/BetafuncEvtVtxGenerator.cc} $z$ is a Gaussian of
  $\sigma_z = 3.6\cm$, $x$ and $y$ follow the $\beta$ function around $\beta^* = 30\cm$, and the time
  is a second, independent Gaussian of the same width, $\sigma_t = 3.6\cm/c = 120\ps$, around
  \code{TimeOffset = 0}.
\item HL-LHC, \code{DBrealisticHLLHC}: \code{HLLHCEvtVtxGenerator}
  (\path{VtxSmearedRealisticHLLHC_cfi.py}, record \code{SimBeamSpotHLLHCObjectsRcd}, payload
  \path{SimBeamSpotHLLHCObjects_z4p3cm_t193ns_14TeV_nominal_mc_v1}: $E_p = 7000\GeV$, crossing angle
  $510\,\mu$rad, crab angle $380\,\mu$rad at 400 MHz, $\beta^* = 0.20\,$m, bunch length $9\cm$)
  computes the overlap density of two crabbed bunches in $(x, y, z, t)$ (\code{integrandCC} in
  \path{HLLHCEvtVtxGenerator.cc}): $\sigma_z = 4.3\cm$ and $\sigma_t = 193\ps$, $z$ and $t$ correlated.
\end{itemize}
This is the number the MTD timing layers of Phase-2 are built against: with a vertex spread of $193\ps$
and a per-track resolution of a few tens of ps, the collision time is what separates the 200 pileup
vertices of a bunch crossing, and the generator stage is where it is drawn. The
non-database block \code{HLLHCVtxSmearingParameters} of \path{VtxSmearedParameters_cfi.py} is still
labelled "for HL-LHC operation at 13TeV" (\code{EprotonInGeV = 6500}); \code{HLLHC14TeV} clones it with $7000$.

\paragraph{Run-2 vintage inside the GEN step.} Of three settings, one is wrong for 2024 and two have not moved since Run 2:
\begin{itemize}
\item \code{LHCTransport} (\code{PPSSimTrackProducer}, added by the \code{ctpps\_2022} modifier of the
  Run-3 era) carries \code{BeamEnergy = 7000}, set in the block
  \code{opticalfunctionsTransportSetup\_2021} of
  \path{SimTransport/PPSProtonTransport/python/OpticalFunctionsConfig_cfi.py}, in a job with
  $6800\GeV$ beams: a 3\% error in the proton optics of the 2024 PPS simulation. For Phase-2 the module
  disappears (\path{Configuration/Eras/python/Era_Phase2_cff.py} excludes \code{ctpps\_2022}); no
  HL-LHC proton transport exists in this release.
\item \code{pythia8CommonSettings} (\path{Pythia8CommonSettings_cfi.py}) keeps
  \code{ParticleDecays:tau0Max = 10} (mm): every hadron with $c\tau > 10\mm$ is left for Geant4 to decay.
  This 2017 convention defines the "stable particle" of all later chapters and has not been revisited
  for the Phase-2 tracker.
\item the generator jets and MET of \code{pgen} have no era dependence: \code{jetPtMin = 3},
  \code{Rho\_EtaMax = 4.5} (\path{GenJetParameters_cfi.py}) and the fiducial cut of \code{genMetCalo},
  $|\eta| < 5.0$ hard-coded in \path{RecoMET/METAlgorithms/src/GenSpecificAlgo.cc}, are the same for
  the 2024 and the Phase-2 detectors ($|\eta| < 4$ tracking, HGCal to $3.0$).
\end{itemize}
Not established from the release: whether the $14\TeV$ fragments will be retuned (there is no Phase-2
tune directory) and the real HL-LHC optics of the first years (the payload is "nominal").
